\begin{table*}[t]
\centering
\caption{Taxonomy of reward hacking modes observed across GRPO training conditions. Each mode produces high prove rate through structurally degenerate formalizations.}
\label{tab:hacking-taxonomy}
\resizebox{\textwidth}{!}{%
\begin{tabular}{llcccp{6cm}}
\toprule
\textbf{Hacking Mode} & \textbf{Experiment} & \textbf{PR} & \textbf{ISR} & \textbf{Comp.} & \textbf{Diagnostic Features} \\
\midrule
Conclusion Embedding & plan001-s42 (exec-only) & 0.93 & 0.1324 & --- & Conclusion literal appears in premises; ISR≈0 for most proved samples (median=0) \\
Tautological Collapse & plan001-s44 / plan013 & 0.98 / 1.00 & 0.5816* / 0.535* & 2.04 / ~2 & 2-component proofs (P$\rightarrow$P); ISR≈0.5 is structural artifact, not genuine quality \\
Repetition / Length Explosion & plan001-s43 (exec-only s43) & 0.14 & 0.0 & --- & Repeated premises; low PR due to parser/solver failure; ISR=0 across all proved \\
Contradiction Explosion & plan016 (1000-step) & 0.96 & 0.0028 & 18.55 & High component count but 94.8% have ISR=0; many redundant contradicting premises \\
Complexity-without-Utility & plan017 (random $\beta$=0.3) & 0.48 & 0.0943 & 9.60 & Many components (9.6 mean) but few necessary; ISR=0 ratio high \\
\bottomrule
\end{tabular}%
}
\vspace{2pt}
\begin{minipage}{\textwidth}
\footnotesize
* ISR values for tautological collapse are artifacts of the 2-component proof structure, not genuine structural quality. \\
\end{minipage}
\end{table*}
